% !TeX root = ../教材题库_全部条目_转写工作稿.tex
% Source: XJ-B1-C05-EX-P160-Y1053
% 第5章 / 5.1 任意角和弧度制 / 教材 p.160 / PDF 第167页
% OCR未校对：公式与文字均可能有误；未完成数学TeX转写；禁止直接发给学生
% 原题及相邻原解截图（核对用）：../题目原图/XJ-B1-C05-EX-P160-Y1053-01.webp
% 本文件只在转写完成后作为题干使用；OCR 条目尚不是可靠题干。
\begin{bankitem}{5.1 例题 例题5}{第5章 / 5.1 任意角和弧度制 / 教材 p.160 / PDF 第167页}{XJ-B1-C05-EX-P160-Y1053}{OCR未校对：公式与文字均可能有误；未完成数学TeX转写；禁止直接发给学生}
\begin{ocrdraft}
[教材p.160 / PDF第167页]
\par
00如图5．1一6，设扇形的圆心角a=x
\par
，半径为0弧
\par
长为/，扇形面积记为S．
\par
（1)用厂与表示扇形的面积S；
\par
（2）用厂与/表示扇形的面积S．
\par
解（1）rad角的大小是周角的
\par
，所形成扇形的面
\par
图5．1
\par
一6
\par
积S是整个圆的面积丌严的一一
\par
1
\par
1
\par
（2）由（1）可得S=
\par
2
\par
2
\par
又因为/=
\par
厂，
\par
所以S=
\end{ocrdraft}
\end{bankitem}
